\chapter{Discussion}
\label{ch:discussion}

\section{The evaluation matters as much as the model}
On all three processes, the largest improvements came from changes to the
evaluation, not to the model. The search exploited every weakness in the
score that it could find: constant flagging, breaking one scenario to win on
average, and fitting the timing of the visible scenarios. On the
quadruple-tank process, a long plateau turned out to be caused by the test
bed itself: faults that were never applied, a decision interval that fixed
a floor on the score, and a base-layer pairing that made one scenario
impossible. A perfect-detection oracle exposed the last two immediately
once it was computed. Checking that the test bed is sound, and bounding what
is achievable, should come before interpreting the progress of the search.

\section{What the LLM needs}
Progress on the quadruple-tank process was much faster after the prompt
gained a plant-engineer-level description of the process and per-decision
traces of the champion's worst scenarios. This agrees with AlphaEvolve and
Eureka, where removing problem context or result feedback degraded the
search \cite{novikov2025alphaevolve,ma2024eureka}. Sampling several
candidates per generation also mattered: the candidates of one generation
often differed in score by a factor of two or more.
\TODO{Quantify this with an ablation of the process description and the traces.}

\section{Reasoning}
Without reasoning, the LLM made steady early progress at about a
twenty-fifth of the cost, but then stalled, wrote physically incorrect
lessons and made avoidable coding mistakes. Because lessons are shown to
later generations as constraints, an incorrect lesson can keep steering the
search in the wrong direction. Low reasoning was enough to reach the oracle
and was about six times cheaper than high reasoning, which also produced
many empty, billed responses.

\section{Interpretability}
Every champion is a single Python function of 43 to 210 lines with named
thresholds and a diagnosis string for each decision. The supervisor that
matched the oracle on the quadruple-tank process
(Appendix~\ref{app:supervisor}) has 43 lines and flags a tank when any of
three conditions on the last 10~s holds: the level is falling while the error
is large; the loop's effort is above the healthy maximum while the error is
large and the level is not recovering; or the pump is close to saturation
and the level is not rising quickly. An engineer can read, test and modify
such a function, and the security sandbox guarantees that it cannot import
modules, touch files or run indefinitely. This is the main advantage over a
learned black-box controller.
\TODO{Discuss how interpretable the longer single- and two-tank supervisors are in practice.}

\section{Limitations}
\begin{itemize}
\item The processes are small simulations with a single fault type (leaks),
  and the quadruple-tank process has noise-free measurements.
\item Each configuration was trained once; the search is stochastic, so
  repeated runs are needed for firm conclusions.
\item The batteries are small (six to eight scenarios each), and the single-
  and two-tank supervisors generalized noticeably worse to the held-out
  scenarios.
\item The tracking error is measured against the active setpoint, so a
  supervisor can reduce it slightly by lowering the setpoint of a leaking
  tank without improving the process. On the quadruple-tank process this is
  worth about 10 points per scenario.
\item No comparison with TD-MPC or any other learned controller has been made
  yet.
\end{itemize}
